\section{研究目标与贡献}

本文以 HuggingFace LeRobot 框架\cite{lerobot}下的 ACT 策略为研究对象，
围绕其在树莓派 5（4 核 ARM Cortex-A76 处理器，4\,GB LPDDR4X，无 GPU）上的在线控制任务，
研究端到端机器人策略在“观测—推理—执行”闭环中的系统级性能瓶颈。
本文关注的不是模型或外设的单点性能，
而是真实控制循环中摄像头、舵机、深度学习推理框架、Python 运行时和 Linux 内核调度如何共同决定系统运行频率。
因此，本文的首要目标是拆分原本混杂在总耗时中的不同来源，
明确各类开销所在层次、触发路径及其对控制闭环的影响。

为实现上述目标，本文构建覆盖 AI 推理框架层、Python 运行时层、OS 内核层与硬件 I/O 层的跨层次性能分析框架。
实验以 SO-101 六自由度机械臂、双目摄像头和树莓派 5 为平台，
复现 LeRobot 在线控制流程，并结合阶段计时、系统调用追踪、内核追踪和必要的内核插桩，
分别分析摄像头异步读取、舵机串口通信、ACT 策略前向推理和主循环等待等环节。
通过分层测量，本文避免仅凭端到端帧率判断瓶颈，
而是从用户态框架、运行时线程、系统调用等待和设备访问路径共同解释性能现象。

在此基础上，本文面向发现的主要阻塞环节设计优化思路，
重点验证异步推理是否能够缓解同步推理对控制循环的阻塞。
优化目标不是改变策略模型结构或依赖更高算力硬件，
而是在保持现有框架和硬件平台基本不变的前提下，提高推理与执行之间的流水线连续性。

本文的主要贡献体现在三个方面：首先，提出覆盖 AI 推理框架层、Python 运行时层、OS 内核层和硬件 I/O 层的四层性能分析模型，
建立面向“观测—推理—执行”紧耦合场景的多工具联合剖析方法。
其次，通过 Python 阶段计时、strace 工具校准、ftrace 追踪与内核自定义插桩，
端到端量化机器人主循环各阶段时间开销，
定位 ResNet18 视觉编码器为第一瓶颈，并揭示 strace 追踪会明显扰动在线控制帧率。
最后，针对同步推理阻塞控制循环的问题设计异步推理优化方案：
实验表明，单次 ACT 完整推理耗时 1924.8\,ms，
其中 ResNet18 视觉编码器占 66.6\%（1281\,ms）；
在 chunk\_size=100 的同步推理配置下实测帧率仅 18.7\,fps，距 30\,fps 目标仍有约 38\% 缺口，
采用异步推理优化后有效帧率提升至 27.42\,fps（相对提升 70\%），基本接近实时目标。

最终，本文形成可复用的边缘机器人系统性能分析流程，
为推理运行时优化、操作系统调度改进和机器人基础软件设计提供实验依据。
